% 第七章：遇到的困难和解决方法
% 负责人：马福泉

\section{遇到的困难和解决方法} \label{sec:difficulties}

\subsection{技术难点}

\subsubsection{缓存行为建模复杂性}
\textbf{困难描述}：现代处理器的缓存行为受到多种因素影响，包括预取、替换策略、缓存一致性协议等，难以准确建模。

\textbf{解决方法}：
\begin{itemize}
    \item 采用基于机器学习的CRP模型，通过大量实验数据训练
    \item 使用分段线性函数简化模型复杂度
    \item 引入在线学习机制，动态调整模型参数
\end{itemize}

\subsubsection{实时性保证}
\textbf{困难描述}：缓存感知调度需要额外的计算开销，可能影响实时性保证。

\textbf{解决方法}：
\begin{itemize}
    \item 优化算法时间复杂度，使用高效的数据结构
    \item 采用预计算和查表技术减少在线计算
    \item 设计增量更新机制，避免重复计算
\end{itemize}

\subsection{工程实现挑战}

\subsubsection{多核同步}
\textbf{困难描述}：多核环境下的数据同步和状态一致性维护困难。

\textbf{解决方法}：
\begin{itemize}
    \item 采用无锁数据结构减少同步开销
    \item 设计分布式状态管理机制
    \item 使用原子操作保证关键数据的一致性
\end{itemize}

\subsubsection{系统集成}
\textbf{困难描述}：调度器需要与操作系统内核紧密集成，涉及底层系统调用和硬件接口。

\textbf{解决方法}：
\begin{itemize}
    \item 开发内核模块形式的调度器
    \item 设计标准化的接口层
    \item 提供用户态测试框架便于调试
\end{itemize}
